% !TeX root = ../papers/nc-p3.tex
\section{標準二文字と自明二文字を持つ全窓}
\label{sec:one-standard-sl2}
$V_i=U_i\oplus S_i$ とし、$S_i$ は同じ $\mathrm{SL}_2$ の標準表現、$U_i$ は二次元自明表現とする。
$J=\left(\begin{smallmatrix}0&1\\-1&0\end{smallmatrix}\right)$ と置く。
四文字テンソルの純 $U$ 成分を $P$ とし、異なる位置 $i,j$ が $U$ である成分を、
残る二位置 $k,l$ での交代形式に $\operatorname{sgn}(i,j,k,l) B_{ij}$ を掛けて表す。
これにより $B_{ji}=-B_{ij}^{\mathsf t}$ となる。
$b_i=B_{i,i+1}$ とし、$p=B_{02},q=B_{13}$ と書く。添字は4を法とする。

\begin{theorem}\label{thm:one-standard-sl2-classification}
この表現型の標準Koszul左右AS正則性は、所定の二隣接切断の表現階数と四つの $b_i$ の可逆性と同値である。
正則窓の添字保持同型類の重複のないパラメータは次の通りで、各類の最小周期は1である。
\begin{enumerate}
\item $p=J+Q$, $Q$ は非零対称行列、$P=0$。
$Q$ の階数1が一類、階数2では $\det Q\in k\setminus\{0,-1\}$ が完全不変量である。
\item $p=J$。追加データは零または射影二元二次形式 $T$ で、$T$ の階数 $0,1,2$ が三類を与える。
階数0は多項式環である。
\item $p=\left(\begin{smallmatrix}0&1\\1&0\end{smallmatrix}\right)$。
追加データは同じく $T=\left(\begin{smallmatrix}a&b\\b&c\end{smallmatrix}\right)$。
$T=0$ は一類である。$b=0$, $T\ne0$ のときは $a,c$ の非零成分数で二類に分かれる。
$b\ne0$ なら $b=1$ として $ac$ と $a,c$ の非零成分数が完全不変量となる。
\end{enumerate}
復元では $b_0=b_1=b_2=J$, $b_3=JM^{\mathsf t}$, $M=-pp^{-\mathsf t}$,
$q=p/\det p$ とする。後二行の $P$ は下記の $T$ に対応するテンソルである。
この表と定理\ref{thm:sl2-window-profiles}の他の表現型が共有する類は多項式環だけである。
\end{theorem}

\begin{proof}[切断と退化境界の排除]
各辺の標準型の関係は二平面 $E_i\subset U_i\oplus U_{i+1}$ と標準表現の積であり、
二つの直和因子は $U_i S_{i+1},S_i U_{i+1}$ に対応する。標準型の切断は、行と列の符号の変更を除き、
\[
 \begin{pmatrix}p&-b_3^{\mathsf t}\\ b_1&q\end{pmatrix},\qquad
 \begin{pmatrix}q&-b_0^{\mathsf t}\\ b_2&-p^{\mathsf t}\end{pmatrix}
\]
であり、各階数は2である。従って各ブロックは相補二辺の $E$ の完全ペアリングを二つの
座標射影で挟んだ行列であり、可逆性はその両射影が全射であることと同値である。
自明型の切断は、$UU$ と交代 $SS$ の順に、右上と左下が非零の境界ベクトル、右下が零、
左上が $P$ の行列である。非零性は命題\ref{prop:one-standard-sl2-slice}による。
階数2の条件は、各 $i=0,1$ で
\[
 P\in kb_i\otimes (U_{i+2}U_{i+3})+(U_i U_{i+1})\otimes kb_{i+2}
 \tag{*}
\]
となる。並びは巡回して読む。

$E_i$ の左・右射影の全射性を $L_i,D_i$ で表す。各頂点で $D_{i-1}$ または $L_i$ が必要である。
両方が欠ければ $u S_i$ と $S_i v$ が連続純関係になる。
各辺でも $L_i$ または $D_i$ が必要である。実際、$b_{i-1},b_i\ne0$ から両射影の階数は少なくとも1である。
両方が1なら $E_i$ はその左右の像直線の直和である。$b_{i-1}$ の右支持はその左像直線に入り、
$b_{i-1}$ 自身が純 $UU$ 関係なので、その終文字は次の $U_i S_{i+1}$ 純関係の始文字にもなる。
いずれも定理\ref{thm:no-consecutive-pure-relations}に反する。

各反対ブロックの可逆性は $L_i L_{i+2}$ と $D_{i-1}D_{i+1}$ のいずれでも表せる（積は論理積）。
$p,q$ が両方特異なら各反対頂点対は、各頂点の必要条件を満たすために、一方で左だけ、他方で右だけを持つ。
全四射影条件を各辺にも一個ずつ配るには各頂点の選択が巡回して一定でなければならず、矛盾する。
両方可逆なら全射影と全境界が可逆となる。一方だけ可逆なら、巡回によりそれを $p$ としてよい。
頂点0,2は両射影を持ち、1,3では反対の側を一個ずつ選ぶ。
$b_i$ の可逆性は $D_{i-1}L_{i+1}$ なので、二つの特異境界は隣接する。

特異辺を0,1に移す。$b_i=JF_i^{\mathsf t}$ で表すと $F_i$ は射影倍率を除き有向四辺の写像として基底変更される。
可逆二辺を恒等に縮めれば、二つの二次元空間を行き来する二本の階数1写像が残る。
各空間で入射像を $ke_0$ にし、出射核を $ke_0$ または $ke_1$ にする。
位置1での一致は連続する純 $b_0,b_1$ を与えるので許されない。
従って $F_1=E_{00}$, $F_2=F_3=I$, $F_0=E_{00}$ または $E_{01}$ としてよい。ここで $E_{ab}=e_a e_b^{\mathsf t}$ である。
両方で $(*)$ の18本の線形条件は階数16で $P=0$ を強制する。
すると各辺の自明型関係の $UU$ 部分は $kb_i$ のみである。
全射でない左射影の像を消す非零 $U_i^*$ でその辺を評価すると、標準型の関係は消え、評価階数が1以下となる。
右も同じで、定理\ref{thm:all-endpoint-ranks-two}により全射影が必要になる。特異境界は全て除かれた。
逆に可逆境界に対する十分性は命題\ref{prop:one-standard-sl2-slice}の全次数の極限論法である。
\end{proof}

\begin{proof}[正規形と周期]
標準文字のスカラーを四つの基底変更に吸収すると、各 $B_{ij}$ には $g_i B_{ij}g_j^{\mathsf t}$、
$P$ には $\bigotimes g_i$ が作用し、さらに窓全体の射影倍率 $\lambda$ がある。
三辺を $J$ に揃え、上記二切断のSchur補を取ると
$q=-Jp^{-1}b_3^{\mathsf t}=-Jp^{-\mathsf t}J=p/\det p$, $M=-pp^{-\mathsf t}$ となる。
残る変更は $g_0=g_2=g$, $g_1=g_3=g/(\lambda\det g)$ である。
$M$ の射影Jordan型に従い $(*)$ を解く。$M=\operatorname{diag}(1,t)$ のとき18行16列の係数行列に
行列式 $-(t-1)^3/t^3$ の小行列があり、非自明単冪型でも階数16となる。
$M$ がスカラーなら核の次元は3で、$\Sym^2 U\otimes\det U$ である。
明示的には $(a,2b,c)$ を、語 $0000,0001,\ldots,1111$ の辞書順で次の三列の係数とする。
各列は表示した語の符号付き和である。
\[
 \begin{aligned}
 &-0001+0010-0100+1000,\\
 &-0101+1010,\\
 &-0111+1011-1101+1110.
 \end{aligned}
\]
$g$ の四重作用をこの三列に掛けると $T\mapsto(\det g)gTg^{\mathsf t}$ を得る。
従って正規化後は $(p,T)\mapsto(\lambda gpg^{\mathsf t},gTg^{\mathsf t}/(\lambda\det g))$ となる。
$g$ のスカラー部分と $\lambda$ により二成分の倍率は独立である。
$M$ がスカラーであることは $p$ が対称または交代であることと同値である。
それ以外は交代部を $J$ にして $p=J+Q$ とすれば、$Q$ には $\mathrm{SL}_2$ の合同が作用する。
二次形式の合同分類と $\det(J+Q)=1+\det Q$ が第一行を与える。
交代の場合は $\PGL_2$ 全体、対称の場合は等方二直線を保つ対角トーラスとその交換が $[T]$ に作用し、表の残りを得る。

一巡回して符号規約の共通符号を吸収すると新しい反対ブロックは $q$ である。
三つ目の境界を $J$ にするには新位置3だけに $M^{-1}$ を作用させればよい。
従って新 $p$ は元の $p$ に比例する。$P\ne0$ のとき $M$ はスカラーであり、上記三列の巡回は全て符号反転である。
独立な倍率を使えば元の記録に戻るので周期は1である。
\end{proof}

\begin{proof}[印を外した同型]
射影安定群のLevi部分群を取り、その半単純導来群の単連結被覆を $D$ とする。
標数零で半単純部分群はLevi内へ共役できる。位置0への射影は $D$ 上有限核である。
実際、正次元の正規核には単純因子の根 $\mathrm{SL}_2$ が含まれ、四位置の型一致に反する。
従って $D$ を位置0の四次元表現で調べられる。
四次元以下に非自明な表現を持つ単純Lie代数はWeyl次元公式により
$\mathfrak{sl}_2,\mathfrak{sl}_3,\mathfrak{sp}_4,\mathfrak{sl}_4$ である。
単純な場合、表の印は一つの共役類になる。$\mathrm{SL}_3$ では三次元標準または双対上で $2+1$、
$\mathrm{SL}_4$ では四次元標準または双対上で $2+1+1$ の表現による。
$\mathrm{Sp}_4$ では標準二平面は自明文字と直交する非退化平面であり、共役で全て移り合う。

複数因子が作用できるのは二つの $\mathrm{SL}_2$ の標準の直和またはテンソル積に限る。
テンソル積ではどの埋込みも表の型にならない。直和では表の印は個別の因子である。
このとき第二因子は四位置とも $U_i$ に標準として作用し、$S_i$ には自明に作用する。
三辺 $J$ により四つの $U_i$ 上の作用は同じになり、$p$ は交代形でなければならない。
$\Sym^2 U$ の不変元は零なので $P=0$ である。このテンソルは四文字の全交代体積で、
多項式環では二つの印は実際に共役である。従ってどの記録も印を外すことで他の記録と一致しない。

斉次包絡環と非自明単冪型について、窓テンソル上の四位置の無限小基底変更の和を零として解くと、
位置0への安定Lie代数の射影次元は順に3,4である（スカラーを除く）。
既に三次元単純因子があり、半単純部分の次元はこれより大きくなれない。
その既知の矢表現はそれぞれ $L_3+L_1,2L_2$ であるから、いずれも本表と重ならない。
全線形計算は \path{computations/verify_one_standard_sl2.py} で再現できる。
\end{proof}

\begin{corollary}\label{cor:semisimple-symmetry-closed-limits}
非自明な連結半単純対称性を持つ全正則窓は半安定であり、軌道閉包の唯一の多重安定類も正則である。
定理\ref{thm:one-standard-sl2-classification}の表では $P$ を零にし、第一行で $\rk Q=1$ なら
さらに $Q$ を零にすればその類を得る。斉次包絡環と非自明単冪型の閉類は多項式環である。
\end{corollary}
\begin{proof}
各位置で $U$ を $t$, $S$ を $t^{-1}$ 倍する特殊線形変換の極限は二標準文字成分を保ち $P$ を消す。
階数1の $Q$ は $\mathrm{SL}_2$ の対角一径数群で零にでき、三つの $J$ と $p$ の交代部を保つ。
残る $p$ は全て非零の二つの反対角成分だけを持つ表示を取れる。非退化対称部の二等方方向を基底にすればよい。
すると各 $b_i,p,q$ も同じ支持で可逆で、四文字テンソルの支持は四文字 $u_0,u_1,s_0,s_1$ の
24通りの順列全体である。全順列の一様平均は四位置全てで均一で、重み多面体の相対内部にある。
従って対角群の軌道はアフィンで閉であり、複素数上では対角変換により各位置の対角モーメントを零にできる。
一位置だけが異なる二語はないから非対角モーメントも零となり、Kempf--Ness判定で全群について閉である。
一般の標数零代数閉体へは有限生成定義体からの基底変更で移る。

斉次包絡環では前節の表示でスカラー文字を $t^3$, 三次元文字を $t^{-1}$ 倍すると
一スカラー成分が残り二スカラー成分が消え、多項式環を得る。
非自明単冪型では重複度の $M$ を行列式1の対角変換で共役して $I$ に落とせる。
非自明半単純群は根部分群から来る $\mathrm{SL}_2$ を含むので二つの定理が全場合を覆う。
述べた極限はすべて非零アフィン閉軌道であり、半安定性と一意性が従う。
\end{proof}
